\documentclass{article}

% if you need to pass options to natbib, use, e.g.:
%     \PassOptionsToPackage{numbers, compress}{natbib}
% before loading neurips_2026

% The authors should use one of these tracks.
% Before accepting by the NeurIPS conference, select one of the options below.
% 0. "default" for submission
\usepackage{neurips_2026}
% the "default" option is equal to the "main" option, which is used for the Main Track with double-blind reviewing.
% 1. "main" option is used for the Main Track
%  \usepackage[main]{neurips_2026}
% 2. "position" option is used for the Position Paper Track
%  \usepackage[position]{neurips_2026}
% 3. "eandd" option is used for the Evaluations & Datasets Track
 % \usepackage[eandd]{neurips_2026}
% 4. "creativeai" option is used for the Creative AI Track
%  \usepackage[creativeai]{neurips_2026}
% 5. "sglblindworkshop" option is used for the Workshop with single-blind reviewing
 % \usepackage[sglblindworkshop]{neurips_2026}
% 6. "dblblindworkshop" option is used for the Workshop with double-blind reviewing
%  \usepackage[dblblindworkshop]{neurips_2026}
% 7. "education" option is used for the Education Track
%  \usepackage[education]{neurips_2026}


% After being accepted, the authors should add "final" behind the track to compile a camera-ready version.
% 1. Main Track
 % \usepackage[main, final]{neurips_2026}
% 2. Position Paper Track
%  \usepackage[position, final]{neurips_2026}
% 3. Evaluations & Datasets Track
 % \usepackage[eandd, final]{neurips_2026}
% 4. Creative AI Track
%  \usepackage[creativeai, final]{neurips_2026}
% 5. Workshop with single-blind reviewing
%  \usepackage[sglblindworkshop, final]{neurips_2026}
% 6. Workshop with double-blind reviewing
%  \usepackage[dblblindworkshop, final]{neurips_2026}
% Note. For the workshop paper template, both \title{} and \workshoptitle{} are required, with the former indicating the paper title shown in the title and the latter indicating the workshop title displayed in the footnote.
% For workshops (5., 6.), the authors should add the name of the workshop, "\workshoptitle" command is used to set the workshop title.
% \workshoptitle{WORKSHOP TITLE}
% 7. Education Track
% \usepackage[education, final]{neurips_2026}

% "preprint" option is used for arXiv or other preprint submissions
 % \usepackage[preprint]{neurips_2026}

% to avoid loading the natbib package, add option nonatbib:
%    \usepackage[nonatbib]{neurips_2026}

\usepackage[utf8]{inputenc} % allow utf-8 input
\usepackage[T1]{fontenc}    % use 8-bit T1 fonts
\usepackage{hyperref}       % hyperlinks
\usepackage{url}            % simple URL typesetting
\usepackage{booktabs}       % professional-quality tables
\usepackage{amsfonts}       % blackboard math symbols
\usepackage{nicefrac}       % compact symbols for 1/2, etc.
\usepackage{microtype}      % microtypography
\usepackage{xcolor}         % colors

\usepackage{float}
\usepackage{amsmath}
\usepackage{graphicx}
\usepackage{capt-of}
\usepackage{longtable}

\newenvironment{tight}
  {\begin{itemize}
    \setlength{\itemsep}{0pt}
    \setlength{\parsep}{0pt}
    \setlength{\topsep}{2pt}
    \setlength{\parskip}{0pt}}
  {\end{itemize}}

% Note. For the workshop paper template, both \title{} and \workshoptitle{} are required, with the former indicating the paper title shown in the title and the latter indicating the workshop title displayed in the footnote. 
\title{Methods of Pos Training Alignment in LLMs}


% The \author macro works with any number of authors. There are two commands
% used to separate the names and addresses of multiple authors: \And and \AND.
%
% Using \And between authors leaves it to LaTeX to determine where to break the
% lines. Using \AND forces a line break at that point. So, if LaTeX puts 3 of 4
% authors names on the first line, and the last on the second line, try using
% \AND instead of \And before the third author name.


\author{%
  Zain Farhan \\
  Department of Computer Science \\
  Lahore University of Management Sciences (LUMS) \\
  \texttt{zainfarhan.work@gmail.com} \\
  % examples of more authors
  % \And
  % Coauthor \\
  % Affiliation \\
  % Address \\
  % \texttt{email} \\
  % \AND
  % Coauthor \\
  % Affiliation \\
  % Address \\
  % \texttt{email} \\
  % \And
  % Coauthor \\
  % Affiliation \\
  % Address \\
  % \texttt{email} \\
  % \And
  % Coauthor \\
  % Affiliation \\
  % Address \\
  % \texttt{email} \\
}


\begin{document}


\maketitle


\begin{abstract}
  In this paper we look at various forms of post training for LLMs and study their effect on our Models.
  We will be doing 3 Main experiments: Direct Preference Optimization (DPO), Proximal Policy Optimization (PPO), Group Relative
  Policy Optimization (GRPO) and then we compare their performances on Safe and Unsafe prompts. We will also see the differences
  between Reinforcement Learning with Verifiable Rewards (RLVR) and Reinforcement Learning from AI
  Feedback (RLAIF).
\end{abstract}

\section{Introduction}
  This is a short study done on the different methods of LLM alignment with 3 main focuses. DPO, PPO, and GRPO. For each method
  we will use collect and present our data with a short analysis of our results. Alongside that we will also explore how much of a differences
  using RLVR vs RLAIF makes.


\section{Methods and Setup}

\subsection{Direct Preference Optimization (DPO)}

\begin{tight}
  \item \textbf{Dataset:} UltraFeedback Preference Pairs (course provided)
  \item \textbf{Model:} Qwen2.5-1.5B-Instruct (Forzen with LoRA)
\end{tight}

\section{Results}

\subsection{Direct Preference Optimization (DPO)}

\begin{table}[t]
    \centering
    \small
    \setlength{\tabcolsep}{4pt}
    \begin{tabular}{l c l c c c c c c}
        \toprule
        Condition & $\beta$ & Budget & Loss & Pref. Acc. &
        KL/Token & RM Score & Length & At Cap \\
        \midrule
        Standard & 0.10 & 1 epoch (91 steps)
        & 0.674 & 0.628 & $1.40{\times}10^{-4}$
        & $1.38 \pm 1.16$ & $155.2 \pm 104.6$ & 43.1\% \\

        Short fork & 0.03 & 581 pairs (37 steps)
        & 0.691 & 0.607 & $-1.02{\times}10^{-4}$
        & $1.31 \pm 1.26$ & $153.7 \pm 105.7$ & 43.8\% \\

        Short fork & 0.10 & 581 pairs (37 steps)
        & 0.686 & 0.624 & $-2.04{\times}10^{-4}$
        & $1.29 \pm 1.21$ & $156.7 \pm 104.2$ & 43.8\% \\

        Short fork & 0.30 & 581 pairs (37 steps)
        & 0.677 & 0.624 & $-4.50{\times}10^{-5}$
        & $1.31 \pm 1.17$ & $151.7 \pm 105.2$ & 42.4\% \\
        \bottomrule
    \end{tabular}
    \caption{DPO $\beta$ ablation. The standard model was trained for
    one epoch. the short forks used the same 581 retained training pairs
    and optimizer budget. KL/token is the sampled-response estimator
    averaged over response tokens. Reward-model score and generated
    length are mean $\pm$ standard deviation over 290 held-out prompts.}
    \label{tab:dpo-beta}
\end{table}

\begin{table}[t]
    \centering
    \small
    \setlength{\tabcolsep}{5pt}
    \begin{tabular}{l c c c c c}
        \toprule
        Model & Preferred Longer & Length Matched &
        Rejected Longer & Overall & DPO Loss \\
        \midrule
        Standard
        & 44.9\% (78) & 57.5\% (80) & 67.1\% (79)
        & 56.5\% (237) & 0.6821 \\

        Length-balanced
        & 43.6\% (78) & 58.8\% (80) & 78.5\% (79)
        & 60.3\% (237) & 0.6760 \\
        \bottomrule
    \end{tabular}
    \caption{Held-out DPO preference accuracy stratified by the relative
    lengths of preferred and rejected responses. Counts are shown in
    parentheses. Nine of the 246 held-out pairs were excluded by the
    prompt-length filter.}
    \label{tab:dpo-length-strata}
\end{table}

We first do a standard full training run for DPO, then we make three short forks for a $\beta$ Study. The data can bee seen in Table\ref{tab:dpo-beta}. We can see
a very similar level of preference accuracy, with the exception of $\beta = 0.03$ which is the lowest at 0.607. This tells us that too low a value will damage performance,
but we cant conclusively say a best value with our data. The Response Lengths are also very similar, however considering that all rows show over 40\% of responses exceed
the length cap (and were hence truncated), this data cant not be meaningfully interpreted.

We then try to determine if the model is being biased with its respond length and learning a false correlation between length and preferred response. We can see in Table\ref{tab:dpo-length-strata}
that overall the model trained on length balanced data performs marginally better than standard. With Length Matched pairs, the accuracy is very close to each other. But for Rejected response being longer,
the length balanced model does much better than standard, even though performance in preferred response longer is quite similar. This shows us that improvement is not uniform across strata, and while
length balancing does help with preference discrimination in, it does not conclusively demonstrate that the standard model learned a length bias

\subsection{Proximal Policy Optimization (PPO)}


\begin{table}[t]
    \centering
    \small
    \setlength{\tabcolsep}{5pt}
    \begin{tabular}{lr}
        \toprule
        Quantity & Value \\
        \midrule
        Wall-clock time & 293.8 s \\
        Peak VRAM (allocated) & 4.91 GB \\
        Mean reward, first 5 / last 5 updates & 1.14 / 0.53 \\
        Mean reward, all 20 updates & $1.24 \pm 1.58$ \\
        Responses at 512-token cap without EOS & 4 / 20 \\
        Maximum clip fraction & 0.000 \\
        Critic explained variance (mean) & $-15.2$ \\
        Critic--return correlation (mean) & 0.089 \\
        \midrule
        Held-out reward: midpoint / PPO & 1.866 / 1.838 \\
        Held-out reward difference (SEM) & $-0.028\ (0.035)$ \\
        Held-out response length: midpoint / PPO & 287.5 / 289.0 \\
        \bottomrule
    \end{tabular}
    \caption{Standard PPO continuation and held-out evaluation on 100
    common prompts. Training statistics cover 20 updates; held-out
    responses use greedy decoding with a 768-token cap.}
    \label{tab:ppo-standard}
\end{table}

\begin{table}[t]
    \centering
    \small
    \setlength{\tabcolsep}{5pt}
    \begin{tabular}{lrrrr}
        \toprule
        Policy & $\epsilon$ & Held-out reward (SEM)
        & $\Delta$Reward (SEM) & Mean length \\
        \midrule
        Standard PPO & 0.20 & $1.838\ (0.125)$ & $-0.028\ (0.035)$ & 289.0 \\
        Short fork & 0.05 & $1.867\ (0.129)$ & $+0.001\ (0.026)$ & 289.3 \\
        Short fork & 0.20 & $1.846\ (0.127)$ & $-0.020\ (0.033)$ & 288.9 \\
        Short fork & 0.50 & $1.847\ (0.129)$ & $-0.019\ (0.030)$ & 289.2 \\
        \bottomrule
    \end{tabular}
    \caption{Held-out evaluation of PPO clipping forks on the same 100
    prompts. $\Delta$Reward is the paired difference against the
    untouched midpoint; SEM denotes standard error of the mean.}
    \label{tab:ppo-clipping}
\end{table}

\begin{table}[t]
    \centering
    \small
    \setlength{\tabcolsep}{5pt}
    \begin{tabular}{rrrrrr}
        \toprule
        $\beta_{\mathrm{KL}}$ & Train reward
        & First 4 / last 4 & Final KL
        & Mean length & EOS rate \\
        \midrule
        0.00 & 0.92 & $1.64 / 0.20$ & 0.0054 & 253.0 & 62\% \\
        0.10 & 1.51 & $1.66 / 1.35$ & 0.0048 & 243.6 & 75\% \\
        0.20 & 1.54 & $1.66 / 1.41$ & 0.0025 & 247.9 & 75\% \\
        \bottomrule
    \end{tabular}
    \caption{KL-pressure forks over eight updates each. Train reward
    is the reward-model score on training rollouts; final KL is the
    reported sampled estimator.}
    \label{tab:ppo-kl}
\end{table}

We see in the standard PPO run in Table\ref{tab:ppo-standard} that our output does not deviate enough to put Clipping into play.
Meaning the output of the new Model does not change enough to warrant clipping. The Critic explained variance is -15.2 which is very poor,
implying that the run did not go well. Compared to our short forks in Table\ref{tab:ppo-clipping} we see that having a lower $\epsilon$ results in
in the only positive SEM reward changes, however the changes are very small and even less apparent for mean length and Held Out reward, no clear winner in clipping.

In the KL pressure experiment, seen in Table\ref{tab:ppo-kl}, we see having higher $\beta_{KL}$ results in having a higher Train Reward, as it forces the model to stay closer to the
original policy as the KL penalty is higher.

\subsection{Group Relative Policy Optimization (GRPO)}

\begin{table}[t]
\centering
\small
\caption{Effect of group size under a fixed budget of 192 cached completions.}
\label{tab:grpo_group_size}
\begin{tabular}{rccccc}
\toprule
$K$ & Informative (\%) & Reward SD & SD Across Partitions
& Baseline SE & No Usable (\%) \\
\midrule
2 & 94.8 & 0.388 & 0.027 & 0.566 & 27.6 \\
4 & 95.8 & 0.551 & 0.017 & 0.400 & 23.2 \\
8 & 95.8 & 0.619 & 0.000 & 0.283 & 16.7 \\
\bottomrule
\end{tabular}
\end{table}

\begin{table}[t]
\centering
\small
\caption{Effect of length normalization on gradient contribution across response lengths.}
\label{tab:grpo_length_gradient}
\begin{tabular}{lcc}
\toprule
Method & Log-log slope of gradient norm vs.\ length
& Gradient share from shorter half (\%) \\
\midrule
Canonical GRPO & $-0.452$ & 83.7 \\
Dr.\ GRPO      & $+0.545$ & 28.6 \\
\bottomrule
\end{tabular}
\end{table}

\begin{table}[t]
\centering
\small
\caption{Held-out evaluation of GRPO length-normalization variants on 100 prompts.}
\label{tab:grpo_norm_eval}
\begin{tabular}{lcccc}
\toprule
Policy & Reward (SEM) & $\Delta$ Reward vs.\ Start (SEM)
& Mean Length & Truncation (\%) \\
\midrule
Start / midpoint & 1.856 (0.124) & --- & 292.9 & 8 \\
Canonical GRPO   & 1.904 (0.120) & $+0.049$ (0.032) & 288.8 & 8 \\
Dr.\ GRPO        & 1.847 (0.127) & $-0.008$ (0.051) & 288.3 & 8 \\
Canonical replicate & 1.844 (0.128) & $-0.012$ (0.042) & 285.4 & 10 \\
\bottomrule
\end{tabular}
\end{table}

We can see in Table\ref{tab:grpo_group_size} that as the Group size increases, Reward SD also increases and "No Usable" (uncertainty of its mean reward falls) decreases, meaning better performance
for us. however since the number of group also decreases with group size, we have fewer groups for training.

In Table\ref{tab:grpo_length_gradient} we see that there is a high bias towards towards the shorter half of our responses (meaning they skew our calculations) in canonical GRPO than Dr. GRPO.
We can also see with the log gradient that for canonical the relationship is negative, meaning the longer response lengths, the less the value of the gradients. Meanwhile, with Dr GRPO it is positive,
resulting in the opposite result.

The truncation of the three models is similar, as can be seen in Table\ref{tab:grpo_norm_eval}, except Canonical Replicate which is at 10\%. Canonical GRPO has higher SEM, the only one which is a positive change from the reference policy,
while Dr GRPO and Canonical both had less then the reference policy.

\subsection{Safety Calibration of the Aligned Policies}

\begin{table}[t]
\centering
\small
\caption{Class-corrected refusal rates on the safety evaluation set.
Intervals are 95\% Wilson intervals.}
\label{tab:t4-corrected}
\begin{tabular}{lcccc}
\toprule
Policy
& Safe refusal (judge)
& Safe refusal (keyword)
& Unsafe non-refusal (judge)
& Unsafe non-refusal (keyword) \\
\midrule
SFT  & 40.4\% [35,47] & 40.4\% & 11.5\% [8,17] & 10.5\% \\
DPO  & 43.6\% [38,50] & 40.0\% & 12.0\% [8,17] & 11.0\% \\
PPO  & 41.2\% [35,47] & 40.4\% & 11.0\% [7,16] & 10.5\% \\
GRPO & 41.6\% [36,48] & 40.0\% & 12.5\% [9,18] & 11.5\% \\
\bottomrule
\end{tabular}
\end{table}

\begin{table}[t]
\centering
\small
\caption{Paired comparisons against SFT on identical prompts.
Rates use the known prompt class and judge-derived labels.}
\label{tab:t4-paired}
\begin{tabular}{llrrrrrr}
\toprule
Policy & Metric & SFT & Policy & Diff.
& Newly flagged & No longer & $p$ \\
\midrule
DPO & Safe refusal       & 0.404 & 0.436 & 0.032 & 13 & 5 & 0.096 \\
DPO & Unsafe non-refusal & 0.115 & 0.120 & 0.005 & 7  & 6 & 1.000 \\
PPO & Safe refusal       & 0.404 & 0.412 & 0.008 & 2  & 0 & 0.500 \\
PPO & Unsafe non-refusal & 0.115 & 0.110 & -0.005 & 0 & 1 & 1.000 \\
GRPO & Safe refusal      & 0.404 & 0.416 & 0.012 & 3  & 0 & 0.250 \\
GRPO & Unsafe non-refusal& 0.115 & 0.125 & 0.010 & 3  & 1 & 0.625 \\
\bottomrule
\end{tabular}
\end{table}

\begin{table}[t]
\centering
\small
\caption{Manual-versus-judge agreement over 240 policy-response pairs.
The distinct-response row excludes duplicate responses.}
\label{tab:t4-audit}
\begin{tabular}{lrrrr}
\toprule
Subset & $n$ & Agreement
& Cohen's $\kappa$ & Ambiguous (manual / judge) \\
\midrule
All policy-response pairs & 240
& 73.8\% [68,79] & 0.523 & 5.4\% / 0.0\% \\
Distinct responses only & 117
& 71.8\% [63,79] & 0.513 & 6.0\% / 0.0\% \\
\bottomrule
\end{tabular}
\end{table}

According to Table\ref{tab:t4-corrected} we can see that there is a massive bias towards refusing responses than accepting them, as Safe refusal
hovers around 40\% and Unsafe non refusal hovers around 11\%. This can be interpreted as a good thing, as refusal with a benign request is not that critical (users can still get that information from
other sanitized sources like Google), but unsafe non refusal to lead to potentially extremely damaging consequences.

In Table\ref{tab:t4-paired} we see both DPO and PPO have P values of 1 for Unsafe Non Refusal and GRPO also has a high p value of 0.625. Meaning that the data is not representative
of a systematic change. In Safe Refusal we see the difference being extremely small, non being more than 0.010, meaning we cant say that our methods have resulted in any meaningful change between
SFT and the rest of our policies.

Table\ref{tab:t4-audit} shows an agreement of upto 70\% with compared to a manual audit, however it should be noted that the AI judge does not use the ambiguous category at all.

\subsection{Feedback Source: RLVR versus RLAIF}

\begin{table}[t]
\centering
\small
\caption{Exact final-answer accuracy, format compliance, and mean response length on GSM8K ($n=300$) and SVAMP ($n=100$).}
\label{tab:t5-acc}
\begin{tabular}{lrrrrrrr}
\toprule
Policy & GSM8K Acc. (SEM) & Format (\%) & Len. &
SVAMP Acc. (SEM) & Format (\%) & Len. & Change (pts) \\
\midrule
SFT   & 27.3\% (0.026) & 34.7 & 275.5 & 44\% (0.050) & 57 & 164.6 & +16.7 \\
RLVR  & 30.3\% (0.027) & 37.7 & 275.5 & 44\% (0.050) & 57 & 163.9 & +13.7 \\
RLAIF & 28.3\% (0.026) & 35.7 & 275.5 & 46\% (0.050) & 59 & 163.8 & +17.7 \\
\bottomrule
\end{tabular}
\end{table}

\begin{table}[t]
\centering
\small
\caption{Paired accuracy comparisons on identical problems using exact McNemar tests.}
\label{tab:t5-mcn}
\begin{tabular}{llrrrrr}
\toprule
Dataset & Comparison & First only correct &
Second only correct & Net & Exact $p$ \\
\midrule
GSM8K & RLVR vs.\ SFT  & 11 & 2 & 9  & 0.0225 \\
GSM8K & RLAIF vs.\ SFT & 6  & 3 & 3  & 0.5078 \\
GSM8K & RLAIF vs.\ RLVR & 3 & 9 & $-6$ & 0.1460 \\
SVAMP & RLVR vs.\ SFT  & 2 & 2 & 0  & 1.0000 \\
SVAMP & RLAIF vs.\ SFT & 3 & 1 & 2  & 0.6250 \\
SVAMP & RLAIF vs.\ RLVR & 2 & 0 & 2  & 0.5000 \\
\bottomrule
\end{tabular}
\end{table}

\begin{table}[t]
\centering
\small
\caption{Controlled response-pair diagnostics comparing the verifier with the AI judge. The clean, correct response is the preferred reference.}
\label{tab:t5-diag}
\begin{tabular}{lrrrrr}
\toprule
Perturbation & Verifier $n$ & Better / Tie / Wrong &
Judge $n$ & Better / Tie / Wrong & Order consistency \\
\midrule
Corrupted reasoning, correct final
& 20 & 0\% / 100\% / 0\%
& 40 & 0\% / 100\% / 0\% & 1.00 \\

Good reasoning, wrong final
& 20 & 100\% / 0\% / 0\%
& 40 & 5\% / 95\% / 0\% & 0.90 \\

Correct answer + persuasive filler
& 20 & 0\% / 100\% / 0\%
& 40 & 0\% / 75\% / 25\% & 1.00 \\

Gold distractor, wrong final
& 20 & 100\% / 0\% / 0\%
& 40 & 20\% / 80\% / 0\% & 0.70 \\
\bottomrule
\end{tabular}
\end{table}

We can see SFT perform the worst while RLVR performs the best in GSM8K in table\ref{tab:t5-acc}, however in we also see RLAIF perform better than RLVR in SVAMP.
In paired accuracy comparison we see that excluding the first comparison, all p values are greater than or equal to 0.5, making the data unreliable and it cannot be used to draw any
conclusive conclusions.

Finally in Table\ref{tab:t5-diag} We can see that corrupted reasoning does not seem to effect, This is expected of Verifier as (depending on the task) it is systematically unable to look deeper, however.
the judge also seems to tie in its responses.For Good reasoning but wrong answer, we see that verifier prefers clean all the time (again, this is expected since the verifier only checks the final answer),
the judge however only prefers clean 5\% of the time while tieing 95\% of the time, which is quite problematic. Correct answer with persuasive reasoning ties in Verifier (again expected), but judges seems to prefer
the perturbation 25\% of the time. Which by itself is not bad, but combined with our previous result it strongly suggests that AI judge may have a strong stylistic bias. Lastly,
Gold distractor, wrong final seems to have 100\% better preference for clean with verifier (again this is expected), but judge Ties 80\% of the time, which shows poor performance.


\subsection{Cross Task Discussion}

\section{Conclusions}

\section{References}

\appendix
\section*{Appendix}

\section{More Method Details}
\label{app:more-method}

\section{Deviations and Caveats}
\label{app:deviations}

\subsection{Direct Preference Optimization (DPO)}

\begin{tight}
  \item We changed the Batch Size to 1 and Gradient Accumulate Steps to 16, keeping effective batch size constant. This was done as our hardware would run out of memory otherwise.
\end{tight}



\end{document}